%% GENERATED by python/paper/build.py from results/ -- do not edit by hand.
%% Rebuild:  .venv\Scripts\python.exe python\paper\build.py --only ArgentinaUniversal_vectorX
%% Source:   results/shocks/{eeOnly,universal}_match_rho1.0000.csv
\begin{table}[!htb]
\centering
\begin{threeparttable}
\caption{Pension system reform, year 2010. (vector $X_i$ calibration)}
\label{table:Argentina:Universal_vectorX}
\begin{tabular}{lccc}
\toprule
\textbf{Scenario} & \textbf{Tax rate} & \textbf{Savings rate} & \textbf{Avg. workweek (hours)} \\
\midrule 
Baseline & 10.92\% & 14.40\% & 42.54 \\[1ex] % row: baseline
Economic Equilibrium & 10.92\% & 14.58\% & 42.50 \\[1ex] % row: reform@ee
Full effect & 12.17\% & 14.12\% & 42.39 \\[1ex] % row: reform@full
\bottomrule
\end{tabular}
\begin{tablenotes}[flushleft]
\footnotesize
\item[] \textit{Note:} The reform permanently shifts $\epsilon = 1-\theta + \theta \eta_1 h_1/h$ at the calibration year, unanticipated. The ``Economic Equilibrium'' scenario solves for the new economic equilibrium path with taxes as in the Baseline scenario, but with the new permanent $\epsilon$; the ``Full effect'' scenario lets taxes be determined by the politico-economic equilibrium. $\rho=1$. The savings rate is savings relative to GDP. Vector-$X_i$ calibration: see the note to Table~\ref{table:Arg:Calib_vectorX}.
\end{tablenotes}
\end{threeparttable}
\end{table}
